\chapter{\MakeUppercase{Results and Discussion}}

\section{Experimental Results}
The Badal system was evaluated against a standard multi-tier AWS deployment. The results demonstrate the system's effectiveness in several areas.

\subsection{Resource Discovery and Mapping}
The system successfully identified all 15 resources in the test environment, including cross-service dependencies. The dependency graph correctly highlighted a "Single Point of Failure" where multiple EC2 instances depended on a single IAM role.

\subsection{Vulnerability Detection}
The system identified 3 critical CVEs associated with an outdated Python runtime in a Lambda function and 2 high-severity vulnerabilities in a Linux package installed on an EC2 instance. These were successfully cross-referenced using the NVD database.

\subsection{AI Remediation Accuracy}
The AI remediation engine provided valid AWS CLI commands for 90\% of the identified issues. For the critical Lambda runtime issue, it correctly suggested the following command:
\begin{lstlisting}
aws lambda update-function-configuration --function-name MyFunction --runtime python3.11
\end{lstlisting}

\section{Comparison with Existing Tools}
\begin{table}[h]
\centering
\caption{Comparison with AWS Native Tools}
\renewcommand{\arraystretch}{1.3}
\begin{tabular}{|l|c|c|c|}
\hline
\textbf{Feature} & \textbf{AWS Config} & \textbf{AWS Inspector} & \textbf{Badal} \\
\hline
Resource Inventory & Yes & No & Yes \\
CVE Scanning & No & Yes & Yes \\
Dependency Graph & No & No & Yes \\
AI Remediation & No & No & Yes \\
Risk Scoring & Basic & Basic & Advanced \\
\hline
\end{tabular}
\end{table}

\section{Discussion}
The integration of graph analysis with vulnerability scanning provides a significantly clearer picture of cloud security. For instance, knowing an EC2 instance has a vulnerability is useful, but knowing that the instance has an IAM role with administrative access to the entire account (visible via the dependency graph) elevates the priority of that fix.

The use of an LLM for remediation is highly effective for common security issues. However, for complex, multi-step architectural changes, the AI's suggestions should be treated as a starting point for expert review. The "advisory-only" nature of the remediation engine is a deliberate design choice to maintain safety.

\section{Limitations}
\begin{itemize}
    \item \textbf{API Rate Limits:} Extensive NVD scanning can be slow due to the public API's rate limits.
    \item \textbf{AWS Region Scope:} Currently, scans are performed on a per-region basis rather than globally.
    \item \textbf{Complexity of Dependencies:} Extremely complex IAM policies (e.g., those using conditions or complex wildcards) can be difficult to map perfectly.
\end{itemize}
